The reduction in option errors after supplying the observed stock path raised a separate question: could a change to the stock forecast improve predictions without using future information? Return-scaling checks and reconstruction of the saved paths found no error in the conversion from annualized log growth to daily prices. Conditioning the initial state on recent returns also did not materially improve five-session central forecasts. Median-based stock MAE changed from \$16.28 to \$16.36 for GS and from \$49.58 to \$49.40 for LLY in that diagnostic (Supplementary Table~\ref{tab:stock_initialization}). The effect of the initial state on the forecast distribution largely disappeared within a few sessions under the fitted transitions. We then compared the retained JumpHMM with unchanged-price, adaptive-volatility, and simple directional benchmarks, using earlier data to select the new models' settings. Validation across 2019--2024 selected an EWMA decay of $0.97$ and the zero-drift limit of the directional regression (Supplementary Table~\ref{tab:stock_selection}). The selected directional and adaptive forecasts were consequently identical. The stock and SPY return features did not justify a directional adjustment under this selection criterion. Testing the frozen models in 2025 supplied 245 five-session origins per ticker, substantially more than the 17 available in the 2026 case study. Across these comparisons, central price errors remained close to the unchanged-price benchmark (Table~\ref{tab:stock_validation}). These results did not establish a repair for the large stock-price misses or a consistent central-forecast advantage from the added predictors.

\begin{table}[H]
\centering
\caption{Five-session stock forecasts in 2025 and August--September 2026. Each ticker has 245 origins in 2025 and 17 in 2026. MAE scores the forecast median; MAE, CRPS, and central 90\% interval width are dollars per share. Scores average dates within each of three 10{,}000-path replicates and then average replicates. Adaptive denotes the Gaussian model with historically selected EWMA decay $0.97$. Historical selection chose zero directional drift, making the selected directional model identical to Adaptive. The unchanged-price benchmark has no uncertainty interval.}\label{tab:stock_validation}
\resizebox{\textwidth}{!}{\begin{tabular}{lllrrrr}
\toprule
Year & Stock & Model & Median MAE & CRPS & Coverage (\%) & Width \\
\midrule
2025 & GS & JumpHMM & 22.31 & 16.22 & 84.4 & 84.71 \\
2025 & GS & Unchanged & 22.61 & 22.61 & -- & -- \\
2025 & GS & Adaptive & 22.62 & 16.32 & 89.8 & 94.86 \\
2025 & LLY & JumpHMM & 39.10 & 29.60 & 69.9 & 100.40 \\
2025 & LLY & Unchanged & 39.37 & 39.37 & -- & -- \\
2025 & LLY & Adaptive & 39.35 & 28.57 & 84.4 & 151.37 \\
2026 & GS & JumpHMM & 16.25 & 13.03 & 100.0 & 129.05 \\
2026 & GS & Unchanged & 15.51 & 15.51 & -- & -- \\
2026 & GS & Adaptive & 15.48 & 15.58 & 100.0 & 176.18 \\
2026 & LLY & JumpHMM & 49.44 & 35.60 & 70.6 & 147.26 \\
2026 & LLY & Unchanged & 49.37 & 49.37 & -- & -- \\
2026 & LLY & Adaptive & 49.31 & 33.81 & 94.1 & 191.38 \\
\bottomrule
\end{tabular}
}
\end{table}

The stock-distribution comparison nevertheless identified a narrower benefit from adaptive volatility. In 2025, LLY's five-session CRPS fell from 29.60 under JumpHMM to 28.57 under the adaptive model, while central 90\% coverage increased from 69.9\% to 84.4\%. Average interval width increased from \$100.40 to \$151.37, so the coverage gain came with a wider range of predicted prices. The lower CRPS indicated that this tradeoff improved the distribution score, although coverage remained below its nominal level. The 2026 comparison showed the same direction of change for LLY: CRPS fell from 35.60 to 33.81 and coverage rose from 70.6\% to 94.1\% (Table~\ref{tab:stock_validation}). The adaptive model had lower date-level CRPS on 131 of 245 LLY origins in 2025 and ten of seventeen in 2026. For GS, however, CRPS increased slightly in 2025, from 16.22 to 16.32, and more substantially in 2026, from 13.03 to 15.58. Both GS models covered every 2026 endpoint, but the adaptive intervals were wider. Full five-session results, including mean-based RMSE and the identical selected directional forecasts, are given in Supplementary Table~\ref{tab:stock_validation_five}. The one-session comparison also showed little separation in central price error from an unchanged-price forecast (Supplementary Table~\ref{tab:stock_validation_one}). The findings therefore supported a ticker-dependent improvement in forecast uncertainty for LLY, not a general improvement in predicting stock-price direction. We retained JumpHMM as the worked example for the IV framework and did not select a replacement stock model from these test outcomes.
\FloatBarrier
